\documentclass[11pt]{article}
\usepackage[margin=1in]{geometry}
\usepackage[T1]{fontenc}
\usepackage{lmodern}
\usepackage{microtype}
\usepackage{booktabs}
\usepackage{tabularx}
\usepackage{xcolor}
\usepackage{fancyvrb}
\usepackage[hidelinks]{hyperref}
\usepackage{natbib}
\newcommand{\TBD}[1]{\textcolor{red}{[TBD: #1]}}

% One paragraph per source line: tools/bibaudit.py splits sentences within a line.

\title{Submit the Audit with the Bibliography:\\Claim-Level Citation Records as a Submission Artifact}
\author{Jacob Moore\\\TBD{affiliation}}
\date{}

\begin{document}
\maketitle

\begin{abstract}
Fabricated references produced by language models have made citation checking an editorial concern, and venues now screen reference lists for works that do not exist. Existence is the easy case. Long before generative AI, roughly one quotation in six in the medical literature misrepresented its source, and a reference that exists but does not say what it is cited for passes every existence check. Automated tools that judge whether a source supports a claim are improving, but an editor running them must first retrieve each source and find the relevant passage, the two steps at which they most often fail. The author has already done both. We propose that authors submit, inside the bibliography file, a record for every citing sentence: the claim made, the passage of the source that supports it, its location, and a verdict. Like a replication package, the record does not certify correctness. It turns verification into a cheap check: whether every citing sentence has a record, whether each quoted passage occurs in the source, and whether the passage supports the claim. We describe a format embedded as comments in BibTeX, release a tool that performs the first two checks deterministically, and apply it to the references of this article.
\end{abstract}

\section{Inaccurate citation predates generative AI}

A citation makes two promises: that the cited work exists, and that it says what the citing sentence attributes to it. The second has long been the weaker. A meta-analysis of 28 studies of medical journal articles estimated that about a quarter of quotations contained an error, about half of them major \citep{jergas2015quotation}, and an update covering 46 studies and some 32,000 quotations put the rate at 16.9\%, with no improvement over time \citep{baethge2025quotation}. In an annotated corpus of 3,063 citations to 100 highly cited biomedical articles, about 39\% contained an accuracy error \citep{sarol2024citationintegrity}. Errors of this kind compound: in one citation network, claims came to be treated as established through citation chains that distorted or omitted the underlying evidence \citep{greenberg2009distortion}.

Language models have added a cruder failure. In 2025 and 2026, references to works that do not exist were found in papers accepted at NeurIPS \citep{ansari2026compound,russinovich2026phantom}, ACL identified more than 100 accepted papers citing non-existent literature during camera-ready checks \citep{acl2026statement}, and ICLR and ICML now treat such references as grounds for desk rejection \citep{iclr2026response,russinovich2026phantom}. Tools that check whether references exist are widely available, though limited by reference extraction, incomplete metadata and database coverage \citep{badalova2026hallucinated}.

Existence checking addresses the visible part of the problem. A model that is asked for a reference supporting a claim can as easily return a real paper that does not support it, and a human author can misremember, overstate or cite a secondary account. These citations resolve to real works with correct metadata, so existence checks pass them.

\section{Verifying support is feasible, but expensive for the reader}

Whether a source supports a claim is a claim-verification problem, studied first on abstracts \citep{wadden2020scifact} and now against full text. Recent systems retrieve the cited work, locate candidate evidence and assign a verdict, typically including a partial-support category \citep{haan2025semanticcite,mocan2026refverifier,sadeghi2026deepsciverify,lee2026reflens,li2026bibagent}. Their reported weaknesses are consistent. Retrieval is incomplete: RefVerifier resolved 57.6\% of references in its test manuscripts to full text \citep{mocan2026refverifier}, and BibAgent infers support for paywalled sources from how other papers cite them rather than from the source itself \citep{li2026bibagent}. Locating the passage is harder in full text than in an abstract \citep{mocan2026refverifier,sadeghi2026deepsciverify}, and the full text is where many claims are supported. Verdicts depend on the verifier: in one evaluation the same outputs yielded unsupported-citation rates from about 3\% to 18\% depending on the verifier, and verifiers agreed far more on what was supported than on what to flag \citep{goo2026citefaithfulness}. On human-annotated biomedical citations, the best model reached 0.59 micro-F1 \citep{sarol2024citationintegrity}.

These are the conditions an editor or reviewer faces when checking a submission: without the sources at hand, without knowledge of which passage the author relied on, and with a verifier whose errors fall disproportionately on the citations worth flagging. The author faces none of the first two. When a sentence is written, the author has the source and knows, or should know, which passage supports it. The information that verification most needs exists at submission time and is then discarded.

\section{Proposal: submit the audit record with the bibliography}

We propose that authors submit, with the manuscript, a claim-level record of every citation, embedded in the bibliography file itself. For each citing sentence the record holds the fields in Table~\ref{tab:fields}. The format is not tied to how the record is produced: an author may write it by hand, or use a model to draft it and check the result.

\begin{table}[t]
\centering
\caption{Fields recorded for each citing sentence.}
\label{tab:fields}
\begin{tabularx}{\linewidth}{@{}lX@{}}
\toprule
Field & Content \\
\midrule
Location & File and line of the citing sentence in the manuscript source \\
Sentence & The citing sentence, verbatim \\
Claim & What the sentence attributes to this source, stated in one line \\
Excerpt & The passage of the source that supports the claim, verbatim, with elisions marked \\
Excerpt location & Page, section, table or figure of the excerpt in the source \\
Verdict & Supported, partially supported, unsupported or contradicted \\
Note & Optional: what a partial verdict lacks, or why the excerpt is indirect \\
Source & Version of the source the excerpt is taken from, per entry \\
\bottomrule
\end{tabularx}
\end{table}

The model is the replication package. Several economics journals require code and data that reproduce the published results, and their data editors check them before acceptance \citep{aea2026datapolicy,es2026reproducibility}; ACM awards badges to papers whose artifacts are evaluated \citep{acm2020badging}. A replication package does not certify that the analysis is right. It makes the claim that the code produces the reported numbers cheap to test, and it makes misreporting a deliberate act rather than an oversight. A citation record does the same for the claim that the source says what the paper says it does.

The record supports three checks, of increasing cost. The first is completeness: every citing sentence in the manuscript has a record, and every record corresponds to a sentence still in the manuscript. This is deterministic and catches records left behind when the text changed. The second is excerpt verification: each quoted excerpt is searched for in the source text. This is also deterministic, needs only the editor's own copy of the source, and catches fabricated or misattributed excerpts, including those a model invented while drafting the record. The third is the support judgment: whether the excerpt supports the claim. This still needs a person or a model, but it is now a local comparison of two short texts rather than a search through a paper, and an editor can sample records rather than read every source.

The record does not by itself establish that a citation is accurate. An excerpt can be genuine and taken out of context, and an author can label an overstatement as supported. What the record changes is that the author's reading of each source is written down where it can be checked, at the level of the individual sentence, and that a misleading record is a statement the author made rather than an omission.

\section{Embedding the record in the bibliography}

We embed the record in the BibTeX file as a comment block above each entry, between fixed markers, so that it travels with the bibliography through every system that already accepts one and is ignored by BibTeX itself. Figure~\ref{fig:example} shows one block from this article's bibliography. The blocks are generated from one JSON file per reference, and a script regenerates them after every edit. Sentences that have moved are shown at their new line; records whose sentence has gone are marked \texttt{STALE}; citing sentences with no record are marked \texttt{NOT AUDITED}. The same script can rebuild the JSON from the comment blocks, so a reader holding only the manuscript and the bibliography can run every check.

\begin{figure}[t]
\centering
\TBD{generated comment block for one entry of this article's references.bib}
\caption{One generated record block, as it appears above its entry in this article's \texttt{references.bib}.}
\label{fig:example}
\end{figure}

Excerpt verification against extracted PDF text cannot rely on exact matching. Text extraction breaks words at line ends, reorders tables, renders small capitals and ligatures inconsistently and drops mathematical notation. Our tool compares word sequences after Unicode normalization and rejoining of hyphenated words, and reports an excerpt as \texttt{MATCH} when it occurs verbatim, \texttt{NEAR} when at least 40\% of its word 4-grams occur in the source, and \texttt{MISSING} otherwise. \TBD{On a separate 104-excerpt audit, genuine excerpts covered 46--100\% of their 4-grams (median 100\%), and the same excerpts checked against the wrong source 0--21\% (median 0\%). Decide whether to report this calibration audit here or recalibrate on this article's own references.} \texttt{NEAR} results are worth a glance: in our experience they are tables, formulas and page breaks, not misquotations.

\section{Case study: the references of this article}

We audited the references of this article with the procedure proposed here. \TBD{Process: who/what drafted records (model, version), what a human checked, time and cost per citation.} \TBD{Results: N citing sentences, K references; verdict counts; excerpt-check counts (MATCH/NEAR/MISSING); drift-check result; what the audit changed in the text.} The bibliography file, with its records, is available as supplementary material, and readers can rerun every check on this article.

\section{Limitations}

Excerpt verification needs the source text, and an editor without access to a source cannot check its excerpts, though publishers usually hold text-mining access to far more of the literature than individual reviewers do. The support judgment remains fallible whether made by a person or a model. The record adds work for authors, although an author who has read each source should be able to supply its supporting passage, and a model can draft records for the author to check. Records drafted by a model share that model's failure modes, which is why the excerpt check is deterministic and the support judgment is left open to the reader. Finally, the comment format is specific to BibTeX; reference managers that export CSL JSON or RIS would need an equivalent field.

\section{Conclusion}

Journals screen for references that do not exist because the check is cheap. Checking whether a real reference supports its claim is expensive for a reader and cheap for the author at the moment of writing. Asking authors to submit that record, in the bibliography file they already submit, moves the work to where it is cheapest and makes the remaining checks quick enough for routine editorial use.

\section*{Declarations}
\TBD{Generative AI use: models used to draft the text, search the literature, draft audit records; human verification. Data and code availability: tools/bibaudit.py, audit records and references.bib; source PDFs not redistributed. Competing interests. Funding.}

\bibliographystyle{plainnat}
\bibliography{references}

\end{document}
